\documentclass[11pt]{article}
\usepackage[margin=1in]{geometry}
\usepackage{xcolor}
\usepackage{hyperref}
\usepackage{enumitem}
\hypersetup{colorlinks=true, linkcolor=black, urlcolor=blue}
\newcommand{\rc}[1]{\medskip\noindent\textit{\textcolor{black!65}{#1}}\par\smallskip}
\newcommand{\ans}{\noindent\textbf{Response.}\ }
\setlength{\parindent}{0pt}
\setlength{\parskip}{4pt}

\title{\vspace{-2em}Response to Reviewers\\[2pt]
\large TNSM-2026-11591, resubmitted as a new manuscript}
\date{}
\begin{document}
\maketitle

We thank the Associate Editor and the three reviewers. The reviews prompted measurements we had
not made, and two of them changed the paper's conclusions rather than its presentation. We say so
plainly here and in the manuscript.

\textbf{The headline claim has changed.} Reviewer 3 asked whether the reported gap survives if a
link cannot carry more than its capacity, and Reviewer 2 asked for a comparison against a strong
baseline at a matched budget. Both questions were fair and neither had been answered. Answering
them shrank the central result and moved the contribution, so this is not a revision of the same
paper: it is the measurement study the old manuscript should have been. The method, the
constellation model and the code are unchanged; the claims built on them are not.

\textbf{Summary of what changed.}
\begin{itemize}[leftmargin=1.4em,itemsep=2pt]
\item A congestion-blind, one-pass baseline was added, with its tolerance swept rather than fixed
(new Section VI-H). It is competitive off distribution, and no prior work in this line reports it.
\item Enforcing capacity feasibility reduces the headline gap from $79.1\%$ measured as delay to
$30.9\%$ measured as delivered rate, stable across three loss models (new Section VI-G).
\item The transfer claim is withdrawn in its original form and replaced by a measured boundary:
decisive in distribution, parity outside it (new Sections VI-H and VI-I).
\item We identify a previously unremarked cause of apparent transfer failure: the decoder's softmin
temperature is not a constant of the method and must scale with the constellation (new Section
VI-I).
\item The nearest prior work, arXiv:2601.21921, is now cited and discussed. It was public five
months before our submission and we had missed it.
\item Thirteen sections merged into eight; a summary table of every headline number opens Section
VI; a claims manifest ties every number in the text to the file that produced it.
\item \emph{What was removed.} The contribution list C1--C4 is gone as a list, and with it three
claims we can no longer support in the form we made them: that the method is the first to formalize
LEO routing as equilibrium amortization (arXiv:2601.21921 precedes it); that inductive transfer to a
$12\times$ larger shell is a result in its own right (a one-pass blind baseline reaches the same
place); and that the reported gap is the operational gap (it is the gap in BPR delay, and under a
hard capacity it is smaller). The claim that the paper distils a general evaluation protocol is
narrowed to the three specific pitfalls we measured. No section was deleted outright; the material
was merged, and every measurement in the previous submission is still in the paper.
\end{itemize}

\hrule
\section*{Reviewer 1}

\rc{1. Clarification on the multipath decode temperature parameter. Is it fixed, tuned on a
validation set, or adapted per scenario? Adding a brief note on the selection rationale and its
sensitivity would improve reproducibility.}
\ans Fixed at $\tau=0.2$ for every experiment, never tuned on a validation set and never adapted
per scenario. That is now stated. The sensitivity half of the question turned out to matter far
more than we expected, so we swept it. The setting that maximises recovery moves by one to two
orders of magnitude between the $132$-satellite shell and the larger ones, and holding it at $0.2$
costs a median of $0.094$, $0.089$ and $0.065$ per instance on three unseen shells. Because that
loss was previously charged to the learned price field, the sweep changed a conclusion and not only
a caption. New Section VI-I; sweep in \texttt{code/experiments/r1\_1\_tau\_sweep.py}.

\rc{2. Discrepancy in reported numbers across tables. The recovered fraction for the 264-satellite
OOD case appears as 93\% in Table VII and 94\% in Tables VI and VIII.}
\ans Each of those tables was computed from its own set of seeds, which the text explained but the
tables did not. A run-set note is now in the caption of each of the three tables, so a reader
comparing them sees the reason without hunting for it. More generally, this class of mismatch is
now checked mechanically: a claims manifest ties every number in the paper to the CSV, column,
filter and aggregation that produce it, and a script recomputes all $65$ of them. Your comment also
made us notice a second, quieter version of the same problem: the recovered fraction was printed as
a percentage in some sections and as a decimal in others. It is now a decimal everywhere, which is
what its name says it is.

\rc{3. Minor inconsistency in the proactive drift crossover statement. Fig.\ 8 does not have axis
labels or units clearly marked. Please ensure the figure axes are self-explanatory and that the
crossover value is visually verifiable.}
\ans The axes did carry labels and units, so we read the comment as being about the second half:
the crossover was asserted in prose and could not be checked against the figure. The crossing point
is now marked, annotated with its value, and the blind reference line is named on the plot. Fig.\ 8.

\rc{4. The broader literature on edge computing, resource pricing, and anti-jamming in LEO/UAV and
RIS-enhanced aerial-ground networks may be relevant [R.1]--[R.4].}
\ans We have expanded Section II considerably, and it now covers the learned-routing neighbourhood
in depth, including the closest prior work we had missed. We have not added the four suggested
references. They concern RIS, SWIPT and anti-jamming rather than routing or traffic engineering,
and citing them where our future-work paragraph mentions integrated aerial-ground deployments would
not help a reader of this paper. We would rather say so than pad the bibliography, and we are happy
to revisit if the Editor sees a connection we have missed.

\rc{5. Limited discussion on scalability of the multipath decode. A brief complexity breakdown for
the decode relative to the GNN forward would be helpful.}
\ans Table II now separates the two: the decode is one Dijkstra per destination, the same class as
the all-or-nothing pass, and the forward is a fixed number of sparse operations independent of the
demand matrix. The per-slot budget is reported in all-or-nothing passes throughout, which is the
unit that dominates at scale, rather than as a single wall-clock number.

\hrule
\section*{Reviewer 2}

\rc{1. Could the authors please improve the organization? There are 13 sections in total.}
\ans Merged to eight. Decomposition folds into System Model; Pitfalls and Setup become a single
Evaluation Protocol section, which is the right shape now that the protocol is part of the
contribution; Deployment, Threats, Discussion and Future Work become one Discussion section.
Nothing was cut.

\rc{2. Could the authors please add a paragraph at the end of Section I to briefly describe the
paper organization?}
\ans Added at the end of Section I.

\rc{3. The related work section is organized in quite specific aspects, which leads to the
learning-based control and quantum walks paragraphs being short and lacking insightful analysis.}
\ans This was the most consequential comment of the three reviews, and we treated it as a signal to
look harder rather than to write more. Doing so turned up arXiv:2601.21921, which trains a graph
network to infer per-link congestion prices from the constellation state in a single forward pass
for LEO mega-constellations: the same core idea as ours, differing in the supervision signal. It
was public five months before our submission and we had not cited it. It is now cited and
discussed. The section is rewritten rather than extended: for a paper whose contribution is a
measurement, related work is the target, so it now states what this line of work consistently
compares against and what it never tests.

\rc{4. Could the authors please further explain the amortization optimization problem and add a
figure to depict the system model if necessary?}
\ans Section III-E now states three properties of the problem that decide what a solution can be:
the expectation is over instances rather than over slots of one deployment, which is where the
method's exposure lies; the constraint is on inference cost rather than total cost; and the object
being optimised is a routing rather than a price field, so decoder error is charged to the model.
That third point is not a formality, and it is what Section VI-I measures. A system model figure
already appears as Fig.\ 2, showing the plane structure, intra- and cross-plane links, and the
drifting hotspot band.

\rc{5. What is the basis for decomposing the amortization optimization problem into sub-problems?
How can we ensure that the final results are consistent?}
\ans Section III-F now argues both halves. The basis is that the map is a composition of steps that
fail for unrelated reasons: a price field can be accurate and decode into a poor routing, the
decode can be faithful and the prediction collapse on an unseen constellation, and all of it can
hold at a fixed demand pattern and give way once the pattern moves. Consistency is enforced
structurally rather than assumed: the sub-problems compose in one direction, every one is scored on
the same end-to-end realized travel time rather than on an intermediate such as price error, and
Section V checks that no composed pipeline ever reports a travel time below the system optimum.

\rc{6. Zero-shot transfer is tested mainly from one 132-satellite shell to larger shells of related
structure. It would be better to add tests across inclination, number of planes, capacity and
demand distributions, and load regimes. The fact that the 53-degree case ``is near static'' weakens
the broader time-varying-topology framing.}
\ans Transfer is now measured across a ladder of shell sizes and across a change of inclination,
and the results changed the paper. We also accept the framing point: the mechanism our method
exploits is time-varying \emph{load}, not time-varying topology, and the manuscript now says so
rather than claiming the broader property. Capacity and demand distributions remain untested and
are stated as a limitation rather than glossed.

\rc{7. One-step correction and geographic greedy are weak comparators. It would be better to
consider stronger iterative-but-budgeted baselines, warm-started MSA, damped or adaptive load
balancing, ECMP-style multipath, and learned models trained directly on routing objectives. It
would also help to compare inference quality at matched computational budgets.}
\ans Agreed, and this comment reshaped the paper. Reading our own cost table as a ladder made the
gap obvious: every cheap comparator we had was weak and every strong one was expensive, with
nothing in the cell our method claims. We added truncated MSA, warm-started MSA, and
tolerance-based multipath, all at matched budgets. Two findings followed. Textbook equal-cost
multipath is a no-op on this topology, since inter-satellite delays take continuous values and
exact ties essentially never occur, so only a tolerance-based variant is meaningful. And with both
sides tuned per shell, the learned field is decisively ahead in distribution ($0.977$ against
$0.668$, paired difference $0.317$) and level outside it ($0.984$/$0.976$, $0.974$/$0.973$,
$0.989$/$0.969$). We report parity rather than advantage off distribution because eight instances
per shell cannot separate those margins from a tie. New Section VI-H.

\hrule
\section*{Reviewer 3}

\rc{1. The Bureau of Public Roads cost function is developed for road traffic and permits offered
load to exceed link capacity. Why is it an appropriate model for packet networks with queues, loss,
transport protocols, and hard link rates?}
\ans It is appropriate for the assignment problem and inappropriate for what we then did with it,
and we should have separated those. On a road, offered load above capacity is a state the system
occupies: vehicles still arrive, later, and the fourth power describes that delay. A packet link
with a hard rate cannot occupy it, because the excess is lost rather than delayed. We now say this
explicitly in Section V and quantify it in Section VI-G rather than defending the choice.

\rc{2. Blind routing is reported to drive some links to nine times their capacity, which produces
the forty-fold headline gap through the fourth-power cost. A data link cannot actually carry nine
times its capacity. You may want to enforce flow feasibility or packet loss and reassess how much
headroom remains.}
\ans We did, and the reviewer is right about the direction. One clarification we would offer: an
offered load of nine times capacity is a real operating state that a moving hotspot produces; what
cannot happen is a link \emph{carrying} it. So the flaw was in the metric rather than in the load
regime. Recomputing the same configurations under a hard capacity, the advantage reads $79.1\%$ as
a reduction in travel time and $30.9\%$ as an increase in delivered rate. Because a loss model is a
choice, we ran three that disagree where it matters: compounding survival along the path gives
$30.9\%$, charging only the bottleneck gives $36.5\%$, and max-min fair sharing gives $37.1\%$. The
effect survives, at well under half its previously reported magnitude, and we report the range
rather than the most favourable member. New Section VI-G.

\rc{3. User equilibrium is repeatedly described as the congestion-aware optimum, although the
system optimum minimizes total travel time.}
\ans Corrected throughout. The system optimum is reported alongside the equilibrium, and it also
now serves a mechanical purpose: every policy in the new experiments is asserted against
$\mathrm{TTT} \ge \mathrm{TTT(SO)}$, which caught a real implementation error during this work.

\rc{4. The proposed features include the load produced by an all-or-nothing routing pass over the
complete demand set. How does this reduce the information requirement relative to solving user
equilibrium, which is criticized for needing the full demand matrix?}
\ans The reviewer is correct and we withdraw that half of the claim. Our features do require the
full demand matrix; what the method removes is the \emph{iteration}, two all-or-nothing passes
against roughly twenty. The manuscript criticised the equilibrium on two grounds and only one of
them is answered by our method, and the introduction now says exactly that.

\rc{5. Several learning-based TE and routing methods are discussed but not included as baselines.
Why are recent methods excluded from the comparison?}
\ans A learned next-hop baseline was already included and remains. Rather than add more learned
comparators, we judged that the missing baseline was the cheap-and-strong one, which is what
Reviewer 2 asked for and what Section VI-H now supplies. We think that is the more informative
addition, and we say in Section VII what it does not cover.

\rc{6. A packet-level validation could be necessary because the main claim concerns routing under
congestion. The current flow-level simulation cannot demonstrate the real performance.}
\ans We agree this is the right eventual test and have not done it; a packet-level study is a
separate piece of work. The hard-capacity models added in Section VI-G are a step in that
direction, not a substitute, and we say so in Section VII rather than implying coverage.

\rc{7. Ground links and gateway bottlenecks are absent even though they can dominate LEO
congestion.}
\ans Also correct, and also out of scope for this manuscript: adding feeder and ground-segment
capacity changes the system model rather than extending it. It is stated as a limitation.

\rc{8. Thirteen main sections make the presentation fragmented.}
\ans Merged to eight, as described under Reviewer 2 comment 1.

\hrule
\section*{Associate Editor}

\rc{The reviewers have recognized the novelty and technical contributions, while they still have
concerns. It is important that TNSM papers contain a comprehensive review of related work.}
\ans We took the related-work sentence literally and it was the most useful line in the letter. It
led us to arXiv:2601.21921, which anticipates our core idea and which we had not cited. Section II
is rewritten around what this literature claims and what it does not test, which is also the frame
of the resubmitted paper.

On the concerns: two of them changed our results rather than our exposition, and the resubmitted
manuscript reports smaller numbers than the original in three places. We think it is a better
paper for it, and we would rather submit the smaller number we can defend.

\end{document}
